\chapter{Conclusion and Future Work}
\label{ch:conclusion}

This thesis presented the design and verification of a deterministic radar-inertial georeferencing coprocessor on the STM32H753 (Arm Cortex-M7). The firmware ingests MAVLink telemetry from the flight controller, hardware-timestampes radar frame-sync pulses, fuses the streams through a 6-state Extended Kalman Filter, converts local polar radar coordinates to global WGS84 coordinates, and streams georeferenced target metadata to the ground station.

\section{Summary of Contributions}

\begin{enumerate}
\item \textbf{A bare-metal georeferencing pipeline} on the STM32H753, implementing DMA-based MAVLink ingestion, timer-synchronized radar timestamping, lock-free inter-stage buffering, EKF sensor fusion, quaternion-based coordinate transformation, and deterministic binary output streaming --- all without an operating system.

\item \textbf{A hardware-in-the-loop verification methodology} using Renode simulation with real MAVLink v2 packets generated by pymavlink, injected via memory-mapped backdoors. This methodology exercises the full pipeline end-to-end with realistic flight data.

\item \textbf{A taxonomy of eleven embedded-systems defects} specific to radar-inertial georeferencing on Cortex-M7, spanning build/toolchain errors, compiler-boundary mismatches, silicon-boundary hazards, and algorithm/protocol defects. Each defect is reproduced, root-caused, and corrected.

\item \textbf{Explicit documentation of simulator-versus-silicon divergences} for the Renode STM32H7 model, with firmware engineered defensively for physical silicon correctness in all cases where the emulator diverges.
\end{enumerate}

\section{Future Work}

\begin{itemize}
\item \textbf{Physical-hardware validation.} Port the firmware to physical STM32H753XI hardware and validate against real Pixhawk MAVLink data and a real radar SoC. The documented Renode-silicon divergences (DMAMUX1, TIM2 input capture, CRC hardware) need physical confirmation.

\item \textbf{Full-state EKF.} Extend to a 15-state EKF estimating position, velocity, attitude, accelerometer bias, and gyroscope bias. This requires quaternion-based EKF formulation (error-state or unscented Kalman filter) to avoid gimbal lock.

\item \textbf{Multi-target tracking.} Extend the output stream to handle multiple radar targets per frame, with data association and track management.

\item \textbf{Formal timing analysis.} Perform worst-case execution time (WCET) analysis on the EKF update step to guarantee the 100\,Hz real-time deadline under all input conditions.

\item \textbf{Integration with ground station.} Implement the ground-station-side binary frame parser and visual overlay to complete the rescue-map use case.
\end{itemize}